\documentclass[12pt,a4paper]{article}

% ─── Pacotes ───────────────────────────────────────────────────────────────────
\usepackage[utf8]{inputenc}
\usepackage[T1]{fontenc}
\usepackage[brazil]{babel}
\usepackage{geometry}
\usepackage{graphicx}
\usepackage{xcolor}
\usepackage{booktabs}
\usepackage{tabularx}
\usepackage{longtable}
\usepackage{array}
\usepackage{enumitem}
\usepackage{titlesec}
\usepackage{fancyhdr}
\usepackage{hyperref}
\usepackage{microtype}
\usepackage{lmodern}
\usepackage{parskip}
\usepackage{amssymb}
\usepackage{mdframed}
\usepackage{tikz}
\usetikzlibrary{shapes.geometric, arrows.meta, positioning}

% ─── Geometria ─────────────────────────────────────────────────────────────────
\geometry{
  top=2.5cm, bottom=2.5cm,
  left=2.8cm, right=2.8cm
}
\setlength{\headheight}{24.02pt}
\addtolength{\topmargin}{-12.02pt}

% ─── Cores ─────────────────────────────────────────────────────────────────────
\definecolor{azulprimario}{RGB}{26, 82, 152}
\definecolor{azulclaro}{RGB}{214, 229, 255}
\definecolor{cinzaclaro}{RGB}{245, 246, 248}
\definecolor{verdeok}{RGB}{39, 174, 96}
\definecolor{amareloaviso}{RGB}{243, 156, 18}
\definecolor{vermelhorrisco}{RGB}{192, 57, 43}
\definecolor{cinzatexto}{RGB}{60, 60, 60}

% ─── Hyperlinks ────────────────────────────────────────────────────────────────
\hypersetup{
  colorlinks=true,
  linkcolor=azulprimario,
  urlcolor=azulprimario,
  citecolor=azulprimario,
  pdftitle={Proposta Técnica — Sistema de Controle de Horas},
  pdfauthor={Manuel Moraes},
}

% ─── Títulos ───────────────────────────────────────────────────────────────────
\titleformat{\section}
  {\Large\bfseries\color{azulprimario}}
  {\thesection.}{0.6em}{}
  [\color{azulprimario}\titlerule]

\titleformat{\subsection}
  {\large\bfseries\color{cinzatexto}}
  {\thesubsection}{0.6em}{}

\titlespacing{\section}{0pt}{1.8em}{0.8em}
\titlespacing{\subsection}{0pt}{1.2em}{0.5em}

% ─── Cabeçalho e Rodapé ────────────────────────────────────────────────────────
\pagestyle{fancy}
\fancyhf{}
\fancyhead[L]{\small\color{azulprimario}\textbf{Sistema de Controle de Horas}}
\fancyhead[R]{\small\color{cinzatexto}Departamento de P\&D  — Confidencial}
\fancyfoot[C]{\small\color{cinzatexto}\thepage}
\renewcommand{\headrulewidth}{0.4pt}
\renewcommand{\headrule}{\hbox to\headwidth{\color{azulprimario}\leaders\hrule height \headrulewidth\hfill}}

% ─── Ambientes auxiliares ──────────────────────────────────────────────────────
\newmdenv[
  backgroundcolor=azulclaro,
  linecolor=azulprimario,
  linewidth=1.2pt,
  leftmargin=0pt, rightmargin=0pt,
  innerleftmargin=12pt, innerrightmargin=12pt,
  innertopmargin=10pt, innerbottommargin=10pt,
  roundcorner=4pt,
  fontcolor=black
]{caixaazul}
\newmdenv[
  backgroundcolor=cinzaclaro,
  linecolor=cinzatexto!40,
  linewidth=0.8pt,
  leftmargin=0pt, rightmargin=0pt,
  innerleftmargin=12pt, innerrightmargin=12pt,
  innertopmargin=10pt, innerbottommargin=10pt,
  roundcorner=4pt,
  fontcolor=black
]{caixacinzanov}

% caixacinza já redefinida acima como caixacinzanov; mantida para compatibilidade
\newmdenv[
  backgroundcolor=cinzaclaro,
  linecolor=cinzatexto!40,
  linewidth=0.8pt,
  leftmargin=0pt, rightmargin=0pt,
  innerleftmargin=12pt, innerrightmargin=12pt,
  innertopmargin=10pt, innerbottommargin=10pt,
  roundcorner=4pt
]{caixacinza}

% ─── Comandos auxiliares ───────────────────────────────────────────────────────
\newcommand{\badge}[2]{%
  \colorbox{#1}{\color{white}\small\textbf{\,#2\,}}%
}

% ══════════════════════════════════════════════════════════════════════════════
\begin{document}

% ─── Capa ──────────────────────────────────────────────────────────────────────
\begin{titlepage}
  \pagecolor{azulprimario}
  \color{white}
  \centering
  \vspace*{3cm}

  {\fontsize{11}{13}\selectfont\uppercase{Proposta Técnica de Desenvolvimento}}\\[0.6em]
  {\fontsize{28}{32}\bfseries\selectfont Sistema de Controle\\[0.2em]de Horas e Atividades}\\[1.2em]
  \color{azulclaro}
  {\large Departamento de TI\enspace|\enspace Envision — setor de Pesquisa e Desenvolvimento}\\[0.4em]
  {\normalsize Versão 1.1\enspace|\enspace Abril de 2026}

  \vfill

  \begin{tikzpicture}
    \draw[white, line width=1pt] (0,0) -- (12,0);
  \end{tikzpicture}

  \vspace{1cm}
  \color{white!80!azulprimario}
  {\small Elaborado por \textbf{Manuel Moraes}\enspace|\enspace Documento confidencial — circulação restrita}
\end{titlepage}
\pagecolor{white}
\color{black}

% ─── Sumário ───────────────────────────────────────────────────────────────────
\tableofcontents
\newpage

% ══════════════════════════════════════════════════════════════════════════════
\section{Sumário Executivo}

Este documento apresenta a proposta técnica para o desenvolvimento de um sistema web
destinado à automação do preenchimento de planilhas de horas e ao rastreamento de atividades
do departamento de P\&D.

O projeto visa eliminar o processo manual baseado em Excel, reduzindo erros, aumentando a
produtividade da equipe e oferecendo ao gestor visibilidade em tempo real sobre a alocação de
pessoas por projeto.

\vspace{0.8em}
\begin{tabularx}{\linewidth}{>{\bfseries}lX}
  \toprule
  Item & Descrição \\
  \midrule
  Empresa      & Envision — setor de Pesquisa e Desenvolvimento  \\
  Usuários     & 11 a 50 colaboradores \\
  Prazo        & 10 semanas (\textasciitilde{}2,5 meses) \\
  Metodologia  & Ágil — Scrum com sprints de 2 semanas \\
  Integração   & Nenhuma integração externa prevista \\
  \bottomrule
\end{tabularx}

% ══════════════════════════════════════════════════════════════════════════════
\section{Problema e Justificativa}

\subsection{Situação Atual}

O controle de horas trabalhadas no departamento de P\&D (internamente denominado \textbf{RD2})
é realizado de forma manual com dois artefatos: um \textbf{Formulário Google Forms} para coleta
de atividades dos membros (de estagiário a gerente) e planilhas Excel para consolidação e
geração de relatórios de custo. Esse processo apresenta as seguintes fragilidades:

\begin{itemize}[leftmargin=1.5em]
  \item \textbf{Retrabalho:} cada colaborador preenche o formulário individualmente, sem
    padronização de formato ou validação de dados.
  \item \textbf{Erros de preenchimento:} ausência de validação automática (horas inválidas,
    projetos inexistentes, datas incorretas).
  \item \textbf{Baixa visibilidade gerencial:} o gestor consolida múltiplas planilhas
    manualmente para obter uma visão consolidada.
  \item \textbf{Planilha de custos manual:} relatórios de custo por projeto são gerados
    a posteriori, com risco de inconsistências.
  \item \textbf{Validação demorada:} o balanço entre os dados do RD2 e os dados do
    Financeiro é feito manualmente — se houver divergência, o processo retorna para revisão.
  \item \textbf{Ausência de rastreabilidade:} não há histórico de alterações ou auditoria
    dos registros.
\end{itemize}

\subsection{Impacto Esperado com a Solução}

\begin{itemize}[leftmargin=1.5em]
  \item Redução significativa do tempo gasto no preenchimento e consolidação de horas.
  \item Eliminação de erros por meio de validações automáticas.
  \item Geração automática da planilha de custos por projeto.
  \item Dashboard em tempo real com alocação de horas por pessoa e projeto.
  \item Histórico completo e auditável de todas as atividades registradas.
\end{itemize}

% ══════════════════════════════════════════════════════════════════════════════
\section{Solução Proposta}

\subsection{Módulos do Sistema}

\begin{tabularx}{\linewidth}{>{\bfseries}lX>{\itshape}l}
  \toprule
  Módulo & Funcionalidades Principais & Perfis \\
  \midrule
  Autenticação (Keycloak) & Login SSO, gerenciamento de perfis e permissões & Todos \\[4pt]
  Lançamento de Horas     & Registro de: projeto, quantidade de horas e justificativa & Colaboradores \\[4pt]
  Gestão de Projetos      & Cadastro de projetos (apenas Gestor), alocação de equipe & Gestores \\[4pt]
  Relatórios / Custos     & Geração automática de planilha de custos, filtros por período & Gestores \\[4pt]
  Dashboard               & Visão consolidada: horas por projeto e colaborador; tabela dinâmica com valor real vs.\ previsto; filtros de período & Gestores \\
  \bottomrule
\end{tabularx}

\vspace{0.6em}
\begin{caixaazul}
  \textbf{Decisão de escopo:} apenas o \textbf{Gestor} tem permissão para cadastrar e editar
  projetos. Os campos de lançamento são \textbf{fixos} (projeto, horas, justificativa) e
  \textbf{não há fluxo de aprovação} — o lançamento é definitivo após a validação automática.
\end{caixaazul}

\subsection{Fluxo Principal}

O sistema substitui o processo atual baseado em Google Forms e Excel pelo fluxo abaixo:

\begin{enumerate}[leftmargin=1.8em]
  \item \textbf{Login via Keycloak} — autenticação com controle de perfis e permissões.
  \item \textbf{Seleção de Projeto} — colaborador visualiza os projetos em que está alocado.
  \item \textbf{Preenchimento de Horas} — registro de: projeto, horas e justificativa
    \textit{(3 campos fixos por projeto)}.
  \item \textbf{Validação Automática} — sistema valida os dados antes de confirmar o lançamento.
  \item \textbf{Geração de Relatório} — planilha de custos gerada automaticamente com base
    nos lançamentos consolidados, substituindo a planilha Excel manual do processo atual.
\end{enumerate}

\textit{Regra especial:} colaboradores alocados em múltiplos projetos visualizarão os 3 campos
separadamente para cada projeto, replicando o comportamento já esperado no formulário atual.

\subsection{Tecnologias Recomendadas}

\begin{tabularx}{\linewidth}{>{\bfseries}lXl}
  \toprule
  Camada & Tecnologia & Justificativa \\
  \midrule
  Frontend       & React.js + TypeScript       & Interface reativa, tipagem segura \\
  Backend        & Node.js + NestJS            & API REST robusta e escalável \\
  Banco de Dados & PostgreSQL                  & Relacional, gratuito e confiável \\
  Autenticação   & Keycloak                    & SSO, conforme especificado \\
  Exportação     & ExcelJS                     & Geração de planilhas \texttt{.xlsx} \\
  Testes         & Jest + Playwright           & Unit, integração e E2E \\
  CI/CD          & GitHub Actions              & Pipeline de lint, testes e deploy \\
  Logging        & Winston / Pino              & Logs estruturados + health checks \\
  Hospedagem     & Docker + VPS ou Cloud       & Escalável, deploy ágil \\
  \bottomrule
\end{tabularx}

\vspace{0.8em}
\textbf{Ambientes:}
\begin{itemize}[leftmargin=1.5em]
  \item \texttt{dev} — local, para desenvolvimento individual.
  \item \texttt{staging} — servidor de homologação; validação pelo cliente antes do deploy.
  \item \texttt{prod} — produção; Docker Compose em VPS como MVP, com opção de migrar para
    AWS/GCP no futuro.
\end{itemize}

% ══════════════════════════════════════════════════════════════════════════════
\section{Justificativa de Decisões Tecnológicas}

\subsection{Critérios de Seleção}

As tecnologias foram escolhidas com base em cinco critérios principais:

\begin{enumerate}[leftmargin=1.8em]
  \item \textbf{Custo Total de Propriedade (TCO):} preferência por ferramentas de código aberto
    ou com licenças gratuitas para uso corporativo.
  \item \textbf{Velocidade de Desenvolvimento:} frameworks e bibliotecas amplamente conhecidas
    no mercado, para reduzir curva de aprendizado.
  \item \textbf{Escalabilidade:} arquitetura preparada para crescimento (até 50 usuários
    simultâneos nesta fase).
  \item \textbf{Manutenibilidade:} código limpo, bem documentado e aderente a boas práticas
    da indústria.
  \item \textbf{Suporte Comunitário:} comunidades ativas, documentação robusta e recursos
    abundantes online.
\end{enumerate}

\subsection{Análise de Alternativas Consideradas}

\begin{tabularx}{\linewidth}{>{\bfseries}l l l l >{\raggedright\arraybackslash}X}
  \toprule
  \textbf{Camada} & \textbf{Escolhido} & \textbf{Alt.~1} & \textbf{Alt.~2} & \textbf{Razão da Escolha} \\
  \midrule
  Frontend   & React.js   & Vue.js  & Angular & Maior demanda de mercado; ecossistema maduro; menos burocracia que Angular \\[4pt]
  Backend    & NestJS     & Express & FastAPI & Arquitetura opinionada; TypeScript nativo; testes integrados desde o início \\[4pt]
  Banco      & PostgreSQL & MySQL   & MongoDB & Relacional~+ JSONB; melhor para auditoria; sem lock-in de schema \\[4pt]
  Autent.    & Keycloak   & Auth0   & Okta    & Open-source; roda on-premises; elimina custo de SaaS por usuário \\[4pt]
  Exportação & ExcelJS    & SheetJS & Aspose  & Documentação excelente; estilos avançados; licença MIT \\
  \bottomrule
\end{tabularx}

\subsection{Decisões Principais}

\begin{caixaazul}
\textbf{1. React + TypeScript (Frontend)}

\textit{Justificativa:} React é o framework mais requisitado no mercado brasileiro e internacional.
TypeScript reduz bugs em tempo de compilação. O ecossistema (Redux, Axios, React Router) é
maduro e bem documentado. Alternativas como Vue.js e Angular foram descartadas porque Vue
tem comunidade menor (poder ficar órfão) e Angular é excessivamente burocrático para um
projeto de 10 semanas.
\end{caixaazul}

\begin{caixaazul}
\textbf{2. NestJS + Node.js (Backend)}

\textit{Justificativa:} NestJS força uma arquitetura limpa (módulos, controllers, services)
sem permitir código ``anárquico'' como Express puro. Suporta TypeScript nativamente. A
estrutura de testes é integrada desde o início. Para Python/FastAPI, teríamos mudança de
stack entre frontend e backend, complicando DevOps (\textit{single language deployment}).
\end{caixaazul}

\begin{caixaazul}
\textbf{3. PostgreSQL (Banco de Dados)}

\textit{Justificativa:} Relacional é correto para este domínio (usuários, projetos, hora\-rios,
auditoria). PostgreSQL oferece JSONB para flexibilidade futura, full-text search e suporta
constraints complexas (garantir que horas $\geq 0$, etc.). MongoDB seria overkill. MySQL fica
para trás em features avançadas.
\end{caixaazul}

\begin{caixaazul}
\textbf{4. Keycloak (Autenticação SSO)}

\textit{Justificativa:} Empresa já mencionou Keycloak na proposta original como requisito.
Keycloak é open-source, roda em Docker, suporta OIDC/OAuth 2.0 padrão. Auth0 e Okta são
SaaS cobrados por usuário/mês e oneram o orçamento conforme o time cresce. Keycloak pode
rodar localmente sem custo adicional.
\end{caixaazul}

\subsection{Diagrama de Arquitetura}

O sistema é organizado em quatro camadas horizontais independentes:

\vspace{1.2em}
\begin{center}
\begin{tikzpicture}[
  client/.style={
    rectangle, rounded corners=3pt,
    draw=azulprimario, fill=azulclaro!30,
    minimum width=2.5cm, minimum height=0.95cm,
    align=center, font=\small\bfseries, inner sep=5pt
  },
  febox/.style={
    rectangle, rounded corners=3pt,
    draw=azulprimario, fill=azulprimario!10,
    minimum width=10.4cm, minimum height=0.95cm,
    align=center, font=\small\bfseries, inner sep=5pt
  },
  authbox/.style={
    rectangle, rounded corners=3pt,
    draw=amareloaviso!80, fill=amareloaviso!20,
    minimum width=3.0cm, minimum height=0.95cm,
    align=center, font=\small\bfseries, inner sep=5pt
  },
  svcbox/.style={
    rectangle, rounded corners=3pt,
    draw=azulprimario, fill=azulclaro!55,
    minimum width=3.0cm, minimum height=0.95cm,
    align=center, font=\small\bfseries, inner sep=5pt
  },
  databox/.style={
    rectangle, rounded corners=3pt,
    draw=verdeok!70, fill=verdeok!15,
    minimum width=2.6cm, minimum height=0.95cm,
    align=center, font=\small\bfseries, inner sep=5pt
  },
  arr/.style={-{Stealth[scale=1.1]}, draw=azulprimario!60, line width=1pt},
  lbl/.style={font=\tiny, fill=white, inner sep=1.5pt, text=cinzatexto},
  layerlbl/.style={font=\scriptsize\itshape, text=cinzatexto!65, anchor=east}
]

%% Camada 1 — Clientes
\node[client] (c1) at (1.5, 4.2)  {Colaborador\\\footnotesize browser};
\node[client] (c2) at (5.2, 4.2)  {Gestor\\\footnotesize browser};
\node[client] (c3) at (8.9, 4.2)  {Admin\\\footnotesize browser};
\node[layerlbl] at (0.15, 4.2) {Clientes};

%% Camada 2 — Frontend
\node[febox] (fe) at (5.2, 2.9)   {React.js + TypeScript\enspace\textbar\enspace Frontend};
\node[layerlbl] at (0.15, 2.9) {Frontend};

%% Camada 3 — Backend
\node[authbox] (kc)  at (1.5, 1.5) {Keycloak\\\footnotesize SSO / OIDC};
\node[svcbox]  (api) at (6.8, 1.5) {NestJS\\\footnotesize API REST};
\node[layerlbl] at (0.15, 1.5) {Backend};

%% Camada 4 — Dados
\node[databox] (db)  at (2.8, 0.0) {PostgreSQL\\\footnotesize banco};
\node[databox] (xl)  at (5.8, 0.0) {ExcelJS\\\footnotesize exportação};
\node[databox] (log) at (8.8, 0.0) {Winston\\\footnotesize logs};
\node[layerlbl] at (0.15, 0.0) {Dados};

%% Setas: Clientes → Frontend
\draw[arr] (c1.south) -- (c1.south |- fe.north);
\draw[arr] (c2.south) -- (fe.north);
\draw[arr] (c3.south) -- (c3.south |- fe.north);

%% Setas: Frontend → Backend
\draw[arr] (fe.south -| kc.north) -- (kc.north)
  node[lbl, pos=0.5, right=2pt] {login};
\draw[arr] (fe.south -| api.north) -- (api.north)
  node[lbl, pos=0.5, right=2pt] {REST};

%% Setas: API → Dados
\draw[arr] (api.south) -- ++(0,-0.3) -| (db.north);
\node[lbl] at (4.8, 1.17) {CRUD};
\draw[arr] (api.south) -- (xl.north)
  node[lbl, pos=0.5, right=2pt] {export};
\draw[arr] (api.south) -- ++(0,-0.3) -| (log.north);
\node[lbl] at (7.9, 1.17) {logs};

\end{tikzpicture}
\end{center}

\vspace{0.4em}
{\small\itshape\color{cinzatexto} Cada camada é independente e se comunica apenas com a camada adjacente. A autenticação é gerida pelo Keycloak antes de qualquer requisição chegar à API.}

\subsection{Fluxo de Dados no Sistema}

O diagrama abaixo representa o fluxo de um lançamento de horas do início ao fim.
As ramificações à direita mostram o que ocorre em caso de erro em cada etapa de validação:

\vspace{1.2em}
\begin{center}
\begin{tikzpicture}[
  stepbox/.style={
    rectangle, rounded corners=4pt,
    draw=azulprimario, fill=azulclaro!35,
    minimum width=8.2cm, minimum height=0.88cm,
    align=center, font=\small\bfseries, inner sep=5pt
  },
  endbox/.style={
    rectangle, rounded corners=4pt,
    draw=verdeok!80, fill=verdeok!18,
    minimum width=8.2cm, minimum height=0.88cm,
    align=center, font=\small\bfseries, inner sep=5pt
  },
  errbox/.style={
    rectangle, rounded corners=3pt,
    draw=vermelhorrisco!70, fill=vermelhorrisco!10,
    minimum width=2.8cm, minimum height=0.88cm,
    align=center, font=\small, inner sep=4pt
  },
  arr/.style={-{Stealth[scale=1.1]}, draw=azulprimario!65, line width=1pt},
  errarr/.style={-{Stealth[scale=1.0]}, draw=vermelhorrisco!70, line width=0.8pt, dashed}
]

\node[stepbox] (s1) at (0, 0.0)
  {\textbf{1.}\ Colaborador preenche formulário\\\scriptsize projeto + horas + justificativa};
\node[stepbox] (s2) at (0, -1.4)
  {\textbf{2.}\ Validação no Frontend\\\scriptsize campos obrigatórios, formato e limites de horas};
\node[stepbox] (s3) at (0, -2.8)
  {\textbf{3.}\ Requisição POST para NestJS API\\\scriptsize token Keycloak no cabeçalho};
\node[stepbox] (s4) at (0, -4.2)
  {\textbf{4.}\ Validação no Backend (NestJS)\\\scriptsize regras de negócio + autorização};
\node[stepbox] (s5) at (0, -5.6)
  {\textbf{5.}\ Persistência no PostgreSQL\\\scriptsize insere registro em \texttt{TimeEntry}};
\node[stepbox] (s6) at (0, -7.0)
  {\textbf{6.}\ Registro em AuditLog\\\scriptsize rastreabilidade imutável da operação};
\node[endbox]  (s7) at (0, -8.4)
  {\textcolor{verdeok}{\checkmark}\ Confirmação exibida ao usuário};

%% Setas do fluxo principal
\foreach \a/\b in {s1/s2, s2/s3, s3/s4, s4/s5, s5/s6, s6/s7}
  \draw[arr] (\a.south) -- (\b.north);

%% Ramificações de erro
\node[errbox] (e2) at (7.2, -1.4)
  {Erro exibido\\\small ao usuário};
\node[errbox] (e4) at (7.2, -4.2)
  {HTTP 400\\\small retornado};

\draw[errarr] (s2.east) -- (e2.west)
  node[midway, above, font=\tiny, text=vermelhorrisco!80] {falha};
\draw[errarr] (s4.east) -- (e4.west)
  node[midway, above, font=\tiny, text=vermelhorrisco!80] {falha};

\end{tikzpicture}
\end{center}

\textit{Nota:} as validações ocorrem em duas camadas (Frontend e Backend) para garantir
tanto a boa experiência ao usuário quanto a segurança no servidor. O AuditLog é gravado
após o commit no banco, garantindo rastreabilidade imutável.

% ══════════════════════════════════════════════════════════════════════════════
\section{Modelo de Dados}

As cinco entidades principais do sistema são descritas a seguir.

\begin{tabularx}{\linewidth}{>{\bfseries\ttfamily}lX}
  \toprule
  Entidade & Responsabilidade \\
  \midrule
  User          & Dados do colaborador; perfil de acesso (admin, gestor, colaborador) \\
  Project       & Projetos cadastrados pelo gestor; status e responsável \\
  ProjectMember & Vínculo entre User e Project; define quais colaboradores estão alocados \\
  TimeEntry     & Registro de lançamento: usuário, projeto, data, horas, justificativa \\
  AuditLog      & Histórico imutável de todas as operações de escrita no sistema \\
  \bottomrule
\end{tabularx}

% ══════════════════════════════════════════════════════════════════════════════
\section{Escopo do Projeto}

\subsection{Dentro do Escopo}

\begin{itemize}[leftmargin=1.5em]
  \item Tela de login integrada ao Keycloak com gerenciamento de perfis.
  \item Cadastro e gerenciamento de projetos (exclusivo para Gestores).
  \item Formulário de lançamento com 3 campos fixos: projeto, horas e justificativa.
  \item Regra de múltiplos projetos: colaboradores em mais de 1 projeto preenchem os campos
    separadamente para cada um.
  \item Dashboard gerencial: horas por projeto e colaborador, com filtros de período.
  \item Geração automática de planilha de custos (Excel) por projeto e período.
  \item Histórico e auditoria de lançamentos.
  \item Controle de permissões por perfil.
  \item Pipeline de CI/CD com ambientes dev, staging e produção.
\end{itemize}

\subsection{Fora do Escopo}

\begin{itemize}[leftmargin=1.5em]
  \item Integração com ferramentas externas (Jira, Azure DevOps, ERP, folha de pagamento etc.).
  \item Aplicativo mobile nativo (iOS/Android) — o sistema será responsivo para uso via browser.
  \item Módulo financeiro avançado (impostos, nota fiscal etc.).
  \item Funcionalidades de inteligência artificial ou análise preditiva.
  \item Etapas manuais de revisão ou reabertura de lançamentos após a validação automática.
  \item \textbf{Validação automática RD2 × Financeiro} — a conciliação entre os lançamentos
    do sistema e os dados da área Financeira é um processo externo e não faz parte do escopo
    desta versão. Poderá ser considerada como evolução futura em uma versão 2.0.
\end{itemize}

% ══════════════════════════════════════════════════════════════════════════════
\section{Requisitos Não-Funcionais}

\begin{tabularx}{\linewidth}{>{\bfseries}l >{\raggedright\arraybackslash}X l}
  \toprule
  Requisito & Critério & Meta \\
  \midrule
  Performance      & Tempo de resposta para lançamento de horas  & $<$ 2 s \\
  Disponibilidade  & Uptime em horário comercial                 & $\geq$ 99,5\% \\
  Concorrência     & Usuários simultâneos sem degradação         & até 50 \\
  Volume           & Lançamentos por mês estimados               & até 5.000 \\
  Compatibilidade  & Navegadores suportados                      & Chrome, Firefox, Edge (últimas 2 versões) \\
  Responsividade   & Uso via browser em dispositivos móveis      & Obrigatório \\
  \bottomrule
\end{tabularx}

% ══════════════════════════════════════════════════════════════════════════════
\section{Conformidade e Segurança}

\subsection{LGPD}

Os dados de colaboradores (nome, e-mail, horas trabalhadas) constituem dados pessoais sob a
Lei Geral de Proteção de Dados (Lei n.º 13.709/2018). O sistema deverá contemplar:

\begin{itemize}[leftmargin=1.5em]
  \item Definição clara de \textbf{controlador} e \textbf{operador} dos dados.
  \item Política de \textbf{retenção}: prazo máximo de armazenamento de registros históricos
    (sugestão: 5 anos, alinhado à política da empresa).
  \item Funcionalidade de \textbf{anonimização ou exclusão} de dados de ex-colaboradores,
    mediante solicitação formal.
  \item Registro de consentimento no aceite de uso do sistema.
\end{itemize}

\subsection{Controle de Acesso}

\begin{tabularx}{\linewidth}{>{\bfseries}lX}
  \toprule
  Perfil & Permissões \\
  \midrule
  Admin       & Gerenciamento de usuários, perfis e configurações gerais do sistema \\
  Gestor      & Cadastro de projetos, alocação de equipe, acesso ao dashboard e relatórios \\
  Colaborador & Lançamento de horas nos projetos em que está alocado; visualização do próprio histórico \\
  \bottomrule
\end{tabularx}

\subsection{Backup e Recuperação}

\begin{itemize}[leftmargin=1.5em]
  \item \textbf{RPO} (Recovery Point Objective): máximo de 24 horas de perda de dados.
  \item \textbf{RTO} (Recovery Time Objective): restauração do serviço em até 4 horas.
  \item Backup diário automático do banco PostgreSQL com retenção de 30 dias.
  \item Testes de restauração a cada sprint de produção.
\end{itemize}

% ══════════════════════════════════════════════════════════════════════════════
\section{Cronograma de Entrega}

O projeto é dividido em \textbf{5 sprints de 2 semanas} (total: 10 semanas), com entrega
validável ao final de cada ciclo.

\vspace{0.5em}
\begin{center}
\begin{tikzpicture}
  % ---- Month header bands (Mes 1: S1-S4, Mes 2: S5-S8, Mes 3: S9-S10) ----
  % Grid left edge x=2.5, week width=1.27, total grid width=12.7
  % Month 1 end: 2.5 + 4*1.27 = 7.58
  % Month 2 end: 2.5 + 8*1.27 = 12.66
  % Month 3 end: 2.5 + 10*1.27 = 15.20
  \fill[azulprimario]      (2.5,  0.45) rectangle (7.58,  0.90);
  \fill[azulprimario!60]   (7.58, 0.45) rectangle (12.66, 0.90);
  \fill[azulprimario!30]   (12.66,0.45) rectangle (15.20, 0.90);
  \node[white,       font=\footnotesize\bfseries] at (5.04,  0.675) {M\^es~1};
  \node[white,       font=\footnotesize\bfseries] at (10.12, 0.675) {M\^es~2};
  \node[azulprimario,font=\footnotesize\bfseries] at (13.93, 0.675) {M\^es~3};

  % ---- Week sub-headers ----
  \foreach \w in {1,...,10} {
    \node[font=\scriptsize, text=cinzatexto]
      at ({2.5+(\w-0.5)*1.27}, 0.22) {S\,\w};
  }

  % ---- Alternating row backgrounds (rows 0, 2, 4) ----
  % Row i: y from -(i*1.1) to -(i*1.1+0.90)
  \fill[gray!8] (0,  0.00) rectangle (15.20, -0.90);
  \fill[gray!8] (0, -2.20) rectangle (15.20, -3.10);
  \fill[gray!8] (0, -4.40) rectangle (15.20, -5.30);

  % ---- Vertical week separators ----
  \foreach \i in {0,1,...,10} {
    \draw[gray!20, very thin] ({2.5+\i*1.27}, 0.45) -- ({2.5+\i*1.27}, -5.35);
  }
  % Stronger month dividers
  \draw[azulprimario!40, thin] (7.58,  0.45) -- (7.58,  -5.35);
  \draw[azulprimario!40, thin] (12.66, 0.45) -- (12.66, -5.35);

  % ---- Sprint labels: name (bold, blue) + phase (small, grey) ----
  % Row centers: -0.45, -1.55, -2.65, -3.75, -4.85
  \node[font=\small\bfseries, text=azulprimario, anchor=west] at (0.05, -0.27) {Sprint~0};
  \node[font=\scriptsize, text=cinzatexto,       anchor=west] at (0.05, -0.65) {Infra \& Setup};

  \node[font=\small\bfseries, text=azulprimario, anchor=west] at (0.05, -1.37) {Sprint~1};
  \node[font=\scriptsize, text=cinzatexto,       anchor=west] at (0.05, -1.75) {Autentica\c{c}\~ao};

  \node[font=\small\bfseries, text=azulprimario, anchor=west] at (0.05, -2.47) {Sprint~2};
  \node[font=\scriptsize, text=cinzatexto,       anchor=west] at (0.05, -2.85) {Lan\c{c}. Horas};

  \node[font=\small\bfseries, text=azulprimario, anchor=west] at (0.05, -3.57) {Sprint~3};
  \node[font=\scriptsize, text=cinzatexto,       anchor=west] at (0.05, -3.95) {Dashboard};

  \node[font=\small\bfseries, text=azulprimario, anchor=west] at (0.05, -4.67) {Sprint~4};
  \node[font=\scriptsize, text=cinzatexto,       anchor=west] at (0.05, -5.05) {Homologa\c{c}\~ao};

  % ---- Gantt bars (each 2 weeks wide, 0.06 inset) ----
  % Sprint 0: x 2.56 -> 5.00, y -0.06 -> -0.84  (center: 3.78, -0.45)
  \fill[azulprimario, rounded corners=2pt] (2.56, -0.06) rectangle (5.00, -0.84);
  \node[font=\scriptsize\bfseries, white] at (3.78, -0.45) {Setup};

  % Sprint 1: x 5.10 -> 7.54  (center: 6.32, -1.55)
  \fill[azulprimario, rounded corners=2pt] (5.10, -1.16) rectangle (7.54, -1.94);
  \node[font=\scriptsize\bfseries, white] at (6.32, -1.55) {Login};

  % Sprint 2: x 7.64 -> 10.08  (center: 8.86, -2.65)
  \fill[azulprimario, rounded corners=2pt] (7.64, -2.26) rectangle (10.08, -3.04);
  \node[font=\scriptsize\bfseries, white] at (8.86, -2.65) {Horas};

  % Sprint 3: x 10.18 -> 12.62  (center: 11.40, -3.75)
  \fill[azulprimario, rounded corners=2pt] (10.18, -3.36) rectangle (12.62, -4.14);
  \node[font=\scriptsize\bfseries, white] at (11.40, -3.75) {Dashboard};

  % Sprint 4: x 12.72 -> 15.16  (center: 13.94, -4.85)
  \fill[azulprimario, rounded corners=2pt] (12.72, -4.46) rectangle (15.16, -5.24);
  \node[font=\scriptsize\bfseries, white] at (13.94, -4.85) {Deploy};

  % ---- Milestone diamonds at right edge of each bar ----
  \node[diamond, fill=verdeok,      draw=white, line width=0.5pt,
        minimum size=8pt, inner sep=0pt] at (5.00,  -0.45) {};
  \node[diamond, fill=verdeok,      draw=white, line width=0.5pt,
        minimum size=8pt, inner sep=0pt] at (7.54,  -1.55) {};
  \node[diamond, fill=verdeok,      draw=white, line width=0.5pt,
        minimum size=8pt, inner sep=0pt] at (10.08, -2.65) {};
  \node[diamond, fill=verdeok,      draw=white, line width=0.5pt,
        minimum size=8pt, inner sep=0pt] at (12.62, -3.75) {};
  \node[diamond, fill=amareloaviso, draw=white, line width=0.5pt,
        minimum size=10pt, inner sep=0pt] at (15.16, -4.85) {};

  % ---- Outer border ----
  \draw[azulprimario!30, thin] (2.5, 0.45) rectangle (15.20, -5.35);
\end{tikzpicture}
\end{center}

\vspace{0.15em}
{\small\textcolor{cinzatexto}{%
  \textcolor{verdeok}{$\diamond$}\enspace Entrega validada\quad
  \textcolor{amareloaviso}{$\diamond$}\enspace Deploy em produ\c{c}\~ao\quad
  S\,= Semana}}
\vspace{0.7em}

\begin{tabularx}{\linewidth}{>{\bfseries}l l >{\raggedright\arraybackslash}X >{\raggedright\arraybackslash}X >{\raggedright\arraybackslash}X}
  \toprule
  \textbf{Sprint} & \textbf{Per\'iodo} & \textbf{Entregas Principais} & \textbf{Crit\'erio de Aceite} & \textbf{Definition of Done} \\
  \midrule
  Sprint~0 & Sem.~1--2 &
    Setup do ambiente, CI/CD, reposit\'orio, Docker, arquitetura, modelagem do banco e
    Keycloak configurado &
    Pipeline verde; Keycloak com perfis criados; schema do banco migrado e revisado &
    Ambientes de desenvolvimento e staging preparados; automações básicas executando; arquitetura e decisões técnicas documentadas \\[6pt]
  Sprint~1 & Sem.~3--4 &
    Login SSO, gest\~ao de perfis, CRUD de projetos e aloca\c{c}\~ao de colaboradores &
    Login funcional em staging; CRUD de projetos aprovado em revis\~ao com o cliente &
    Controle de acesso por perfil aplicado; cadastro de projetos restrito a gestor; auditoria básica de autenticação e escrita funcionando \\[6pt]
  Sprint~2 & Sem.~5--6 &
    Lan\c{c}amento de horas, valida\c{c}\~oes e regra de m\'ultiplos projetos &
    Formul\'ario aprovado pelo usu\'ario piloto; cobertura de testes $\geq$ 80\% &
    Fluxo com campos fixos entregue; validações automáticas cobrindo cenários principais; múltiplos projetos tratados em blocos separados \\[6pt]
  Sprint~3 & Sem.~7--8 &
    Dashboard gerencial, exporta\c{c}\~ao Excel, relat\'orios por filtro
    (projeto, colaborador, per\'iodo) &
    Dashboard exibindo KPIs reais; planilha validada pelo Financeiro\,/\,RD2 &
    KPIs e filtros homologados; exportação aderente ao template confirmado; logs estruturados e health checks disponíveis \\[6pt]
  Sprint~4 & Sem.~9--10 &
    Testes de aceita\c{c}\~ao (UAT), corre\c{c}\~oes, documenta\c{c}\~ao e
    deploy em produ\c{c}\~ao &
    UAT aprovado; sistema em produ\c{c}\~ao; documenta\c{c}\~ao de opera\c{c}\~ao entregue &
    Operação documentada; backup e restauração verificados; plano de rollback conhecido pela equipe \\
  \bottomrule
\end{tabularx}

\textbf{Prazo total estimado: 10 semanas (\textasciitilde{}2,5 meses)}, dentro da janela de 1 a 3 meses
solicitada.

% ══════════════════════════════════════════════════════════════════════════════

\section{Riscos e Mitigações}

\begin{tabularx}{\linewidth}{Xllp{5cm}}
  \toprule
  Risco & Prob. & Impacto & Mitigação \\
  \midrule
  Escopo mal definido / mudanças frequentes &
    \badge{amareloaviso}{Média} & Alto &
    Reunião de alinhamento antes de cada sprint. Mudanças entram no backlog da próxima sprint. \\[6pt]
  Baixa adoção pelos usuários &
    \badge{amareloaviso}{Média} & Médio &
    Treinamento de 1h + manual de uso. Interface simplificada e intuitiva. \\[6pt]
  Complexidade de configuração do Keycloak &
    \badge{verdeok}{Baixa} & Alto &
    Sprint 0 dedicado à configuração e validação antes do desenvolvimento principal. \\[6pt]
  Indisponibilidade da equipe &
    \badge{verdeok}{Baixa} & Médio &
    Cronograma com buffer absorvido no Sprint 4. \\[6pt]
  Vazamento de dados de colaboradores (LGPD) &
    \badge{verdeok}{Baixa} & Alto &
    HTTPS obrigatório; autenticação via Keycloak; auditoria de acessos habilitada. \\
  \bottomrule
\end{tabularx}

% ══════════════════════════════════════════════════════════════════════════════
\section{Pontos em Aberto para Validação}

Os pontos abaixo permanecem em aberto e devem ser respondidos antes do início das frentes
que dependem diretamente dessas decisões.

\begin{tabularx}{\linewidth}{clXX}
  \toprule
  \# & Status & Questão & Impacto \\
  \midrule
  2 & \badge{amareloaviso}{Aberto} &
    A planilha de custos existente será o template de saída do sistema? &
    Impacta diretamente o Sprint~3, a implementação com ExcelJS e a validação com RD2/Financeiro \\[6pt]
  5 & \badge{amareloaviso}{Aberto} &
    O sistema atenderá apenas este departamento ou múltiplas equipes/empresas? &
    Impacta arquitetura de banco, autenticação, segregação de dados e decisões do Sprint~0 \\
  \bottomrule
\end{tabularx}

\vspace{0.6em}
\begin{caixacinza}
  \textbf{Decisões incorporadas ao escopo (pontos 1, 3 e 4 já respondidos):}\\[6pt]
  \textbf{1.} Apenas o Gestor pode cadastrar projetos.\\[2pt]
  \textbf{3.} Os campos de lançamento são fixos: projeto, horas e justificativa.\\[2pt]
  \textbf{4.} Não há fluxo de aprovação — o lançamento é definitivo após validação.
\end{caixacinza}

\vspace{0.6em}
\begin{caixacinza}
  \textbf{Decisões pendentes que precisam de validação formal:}\\[6pt]
  \textbf{2.} Confirmar se a planilha de custos atual será reutilizada como template oficial de saída.\\[2pt]
  \textbf{5.} Confirmar se o sistema seguirá arquitetura single-tenant ou se precisará suportar múltiplas equipes/empresas.
\end{caixacinza}

% ══════════════════════════════════════════════════════════════════════════════
\section{Próximos Passos}

\begin{enumerate}[leftmargin=1.8em]
  \item Validar e aprovar esta proposta com os stakeholders.
  \item Confirmar disponibilidade da equipe de desenvolvimento.
  \item Validar com RD2/Financeiro se a planilha de custos atual será o template oficial do Sprint~3.
  \item Confirmar com TI/Arquitetura se o produto será single-tenant ou preparado para múltiplas equipes/empresas.
  \item Coletar o arquivo da planilha de custos atual, caso a decisão do item anterior seja positiva.
  \item Definir ambiente de homologação e credenciais do Keycloak.
  \item Agendar reunião de kickoff para o Sprint~0 (ambiente + arquitetura).
\end{enumerate}

\vspace{1em}
\begin{caixaazul}
  \centering
  \textbf{Após a aprovação desta proposta e a resolução dos dois pontos em aberto,\\
  o projeto poderá ter início imediato.}
\end{caixaazul}

\vspace{2em}
\noindent\rule{\linewidth}{0.4pt}\\[0.4em]
{\small Proposta elaborada com base no fluxo do Projeto XL\enspace|\enspace
Abril de 2026\enspace|\enspace\textit{Confidencial — circulação restrita}}

\end{document}
